\documentclass[10pt]{article}
\providecommand{\FIGDIR}{fig}
\newcommand{\DIGESTDATE}{19 September 2026}
\newcommand{\PERIODSTART}{13 September 2026}
\newcommand{\PERIODEND}{19 September 2026}
% ===================== PM Research Watch - shared preamble =====================
\usepackage[T1]{fontenc}
\usepackage[utf8]{inputenc}
\usepackage{amsmath,amssymb}
\usepackage{tgpagella}
\usepackage{tgheros}
\renewcommand{\sfdefault}{qhv}
\renewcommand{\ttdefault}{lmtt}
\usepackage{microtype}
\usepackage[english]{babel}

\usepackage[a4paper,top=17mm,bottom=17mm,left=17mm,right=17mm,
            headsep=6mm,footskip=10mm]{geometry}
\usepackage{graphicx}
\usepackage{xcolor}
\usepackage{colortbl}
\usepackage{tikz}
\usetikzlibrary{calc}
\usepackage[most]{tcolorbox}
\usepackage{titlesec}
\usepackage{fancyhdr}
\usepackage{booktabs}
\usepackage{array}
\usepackage{longtable}
\usepackage{enumitem}
\usepackage{multicol}
\usepackage{ragged2e}
\usepackage[hidelinks]{hyperref}

% ------------------------------------------------------------------ palette
\definecolor{Ink}    {HTML}{14262E}
\definecolor{Deep}   {HTML}{0F4C5C}
\definecolor{Teal}   {HTML}{1F7A8C}
\definecolor{Sky}    {HTML}{4FA3C4}
\definecolor{Sage}   {HTML}{6FA287}
\definecolor{Amber}  {HTML}{E3A008}
\definecolor{Clay}   {HTML}{B5643C}
\definecolor{Coral}  {HTML}{D1495B}
\definecolor{Violet} {HTML}{6A5B9E}
\definecolor{Slate}  {HTML}{4A6572}
\definecolor{Paper}  {HTML}{FBF9F5}
\definecolor{Mist}   {HTML}{EEEAE2}
\definecolor{Rule}   {HTML}{D9D3C9}

\pagecolor{Paper}
\color{Ink}

\hypersetup{colorlinks=true, linkcolor=Deep, urlcolor=Teal, citecolor=Teal,
            pdftitle={PM Research Watch}, pdfauthor={Automated daily digest}}

% ------------------------------------------------------------------ headings
\titleformat{\section}
  {\sffamily\bfseries\large\color{Deep}}{}{0pt}
  {}[{\vspace{1.2mm}\color{Amber}\titlerule[1.1pt]}]
\titlespacing*{\section}{0pt}{5.5mm}{2.8mm}

\titleformat{\subsection}
  {\sffamily\bfseries\normalsize\color{Teal}}{}{0pt}{}
\titlespacing*{\subsection}{0pt}{3.6mm}{1.4mm}

\setlist[itemize]{leftmargin=4.2mm, itemsep=1.1mm, topsep=1.3mm, parsep=0pt,
                  label=\textcolor{Teal}{\footnotesize$\blacktriangleright$}}

% ------------------------------------------------------------------ header/footer
\pagestyle{fancy}
\fancyhf{}
\renewcommand{\headrulewidth}{0pt}
\renewcommand{\footrulewidth}{0pt}
\fancyhead[L]{\sffamily\scriptsize\color{Slate}PM RESEARCH WATCH\;\textcolor{Rule}{\textbar}\;\DIGESTDATE}
\fancyhead[R]{\sffamily\scriptsize\color{Slate}Particulate matter monitoring \& health effects}
\fancyfoot[C]{\sffamily\scriptsize\color{Slate}\thepage}

% ------------------------------------------------------------------ components
% metric tile
\newcommand{\metric}[3]{%
  \begin{tcolorbox}[enhanced, colback=white, colframe=Rule, boxrule=0.5pt,
    arc=1.6mm, left=2.2mm, right=2.2mm, top=1.8mm, bottom=1.8mm,
    borderline west={1.7mm}{0pt}{#3}, nobeforeafter, valign=center,
    height=17mm, width=\linewidth]
    \RaggedRight
    {\sffamily\bfseries\fontsize{17}{17}\selectfont\color{#3}#1}\\[0.5mm]
    {\sffamily\fontsize{6.5}{7.4}\selectfont\color{Slate}#2}
  \end{tcolorbox}}

% pull-quote / key insight box
\newtcolorbox{keybox}[1][]{enhanced, colback=Mist, colframe=Mist,
  arc=1.6mm, left=3.2mm, right=3.2mm, top=2.4mm, bottom=2.4mm,
  borderline west={1.7mm}{0pt}{Amber}, #1}

% subtopic band
\newcommand{\band}[2]{%
  \vspace{3.2mm}\noindent
  \begin{tcolorbox}[enhanced, colback=#2, colframe=#2, arc=1.2mm,
    left=2.6mm, right=2.6mm, top=1.1mm, bottom=1.1mm, boxrule=0pt,
    width=\linewidth, nobeforeafter]
    {\sffamily\bfseries\small\color{white}#1}
  \end{tcolorbox}\vspace{1.6mm}}

% one paper entry: {short cite}{title}{journal}{doi}{body}
\newcommand{\paperentry}[5]{%
  \par\vspace{1.9mm}\noindent
  {\sffamily\bfseries\footnotesize\color{Ink}#1}\;%
  {\sffamily\fontsize{7}{8}\selectfont\color{Slate}#3}\\[0.4mm]
  {\footnotesize\itshape\color{Deep}#2}\\[0.6mm]
  {\footnotesize #4}\\[0.3mm]
  {\fontsize{6.6}{7.6}\selectfont\color{Slate}\texttt{doi:}\,\href{https://doi.org/#5}{\texttt{\detokenize{#5}}}}
  \par}

% caption style
\newcommand{\figcap}[1]{\par\vspace{1.0mm}
  {\fontsize{7.2}{8.6}\selectfont\color{Slate}\textsf{\textbf{Fig.}}\; #1\par}\vspace{1.5mm}}

\providecommand{\permil}{\textperthousand}

\setlength{\parindent}{0pt}
\setlength{\parskip}{0pt}
\renewcommand{\arraystretch}{1.18}


\begin{document}

% ==================================================================== MASTHEAD
\thispagestyle{empty}
\begin{tikzpicture}[remember picture, overlay]
  \fill[Deep] (current page.north west) rectangle
      ($(current page.north east)+(0,-67mm)$);
  \fill[Teal] ($(current page.north west)+(0,-67mm)$) rectangle
      ($(current page.north east)+(0,-68.4mm)$);
  \foreach \x/\y/\r in {22/12/0.9, 41/26/1.5, 63/9/0.7, 88/31/1.1, 112/17/1.9,
                        138/29/0.8, 159/13/1.3, 178/54/1.0, 196/21/1.6,
                        30/44/0.6, 74/57/1.2, 125/48/0.9, 168/6/0.8, 191/59/0.7,
                        52/18/1.0, 100/60/0.6, 148/20/1.1, 66/49/1.4, 205/45/0.8}
    \fill[Sky, opacity=0.30] ($(current page.north west)+(\x mm,-\y mm)$) circle (\r mm);
\end{tikzpicture}

\vspace*{4mm}
{\sffamily\fontsize{7.6}{9}\selectfont\color{Sky}%
 WEEKLY ROLLUP\;\textbullet\;PhD RESEARCH WATCH\par}
\vspace{2.2mm}
{\sffamily\bfseries\fontsize{27}{29}\selectfont\color{white}%
 Particulate Matter\par}
{\sffamily\fontsize{27}{29}\selectfont\color{Sky}%
 Week in Review\par}
\vspace{5.0mm}
{\sffamily\fontsize{8.4}{10}\selectfont\color{white}%
 \mbox{Week ending \DIGESTDATE}\;\;\textcolor{Teal}{\textbar}\;\;\mbox{\PERIODSTART{} to \PERIODEND}\;\;
 \textcolor{Teal}{\textbar}\;\;\mbox{148 records pooled from 7 daily issues}\par}
\vspace{18mm}

% -------------------------------------------------------------- metric strip
\noindent
\begin{minipage}[t]{0.190\linewidth}\metric{148}{RECORDS POOLED\\(148 UNIQUE DOIs)}{Deep}\end{minipage}\hfill
\begin{minipage}[t]{0.190\linewidth}\metric{44\%}{MEASUREMENT SIDE\\(65 RECORDS)}{Teal}\end{minipage}\hfill
\begin{minipage}[t]{0.190\linewidth}\metric{10}{SUBTOPICS\\ACTIVE}{Sky}\end{minipage}\hfill
\begin{minipage}[t]{0.190\linewidth}\metric{47}{POOLED EFFECT\\ESTIMATES (42 WITH A CI)}{Amber}\end{minipage}\hfill
\begin{minipage}[t]{0.190\linewidth}\metric{15}{ABSTRACT-FREE\\(10\% OF WEEK)}{Clay}\end{minipage}

\vspace{6mm}

% -------------------------------------------------------------- week in one box
\section*{\sffamily\bfseries\color{Deep}The week in one paragraph}
\vspace{-1.5mm}{\color{Teal}\rule{\linewidth}{1.1pt}}\vspace{2.5mm}

\begin{keybox}
{\sffamily\bfseries\small\color{Deep}Last week's two open questions both resolved, and
both in the same direction: the collapse in exposure modelling was an artefact, and the
abstract-free problem was fixable. \textit{Exposure} went 20 $\rightarrow$ 38 and
\textsc{metadata only} went 16\% $\rightarrow$ 10\%. The pipeline change that did the
second is the more consequential of the two.}\\[1.4mm]
\footnotesize
\textbf{Scale.} \textbf{148 records} over seven issues on daily totals of 21, 16, 22, 21,
18, \textbf{30} and 20 --- a 1.9$\times$ range, the flattest week this watch has recorded
and flatter again than last week's 2.6$\times$. Against recent rollups: 148 after 126,
150, 109 and 112, so the level is high-normal. The flatness matters more than the level,
because every cluster movement below is being read against a stable denominator rather
than through a single deposit spike. All \textbf{ten} subtopics are populated.

\textbf{\textit{Exposure assessment \& modelling} rebounded 20 $\rightarrow$ 38 and
overshot.} The 12 September rollup recorded a halving from 47 to 20, declined to call it
a turn in the field, and set an explicit test: whether \textit{Exposure} returns to the
mid-20s. It did not return to the mid-20s --- it went past the pre-collapse level to
\textbf{38}, the second-highest reading in the series. \textbf{The test is answered and
the answer is unambiguous}: the 20 was a one-week sampling floor, not a contraction.
\textit{Sensing} moved only 31 $\rightarrow$ 27, so the two measurement clusters again
moved in opposite directions, and the measurement-side share reads 47, 54, 40,
\textbf{44\%} across four weeks --- still oscillating, still not trending, and the 29
August warning against building a claim on it still stands.

\textbf{The abstract-free problem was halved by a pipeline change, not by the
publishers.} \textbf{15 of 148 (10\%)} shipped as \textsc{metadata only}, against 20 of
126 (16\%) last week, and inside \textit{Sensing} the count fell \textbf{10
$\rightarrow$ 4}. Last week's rollup named an abstract-retry pass as the highest-value
outstanding pipeline work. What actually delivered the improvement was something else:
the PMC and AGRICOLA retro-index gate moved into the screening script on 18 September, so
records that used to enter abstract-less and be carried are now rejected at source as
out-of-window artefacts. \textbf{That is a different mechanism with the same visible
effect}, and the distinction matters --- the retry pass is still not built, and the five
same-day Elsevier, ACS and Springer deposits in the 19 September issue confirm the
underlying publisher behaviour is unchanged.

\textbf{Estimate density returned to the series norm and precision improved sharply.}
\textbf{47 pooled estimates against 148 records}, one per 3.1, from a four-week series of
2.9, 4.4, 2.7, \textbf{3.1}. More useful than the density: \textbf{42 of 47 carry a
confidence interval (89\%)}, against 40 of 47 last week and far above the series floor.
\textbf{Four of the five bare points come from a single paper} --- the wildfire versus
non-wildfire PM$_{2.5}$ hospitalisation analysis of 18 September, whose abstract reports
eight headline risks and not one interval.
\end{keybox}

\clearpage

% -------------------------------------------------------------- trend
\section*{\sffamily\bfseries\color{Deep}Subtopic trend and inflection}
\vspace{-1.5mm}{\color{Teal}\rule{\linewidth}{1.1pt}}\vspace{2.5mm}

\noindent\includegraphics[width=\linewidth]{\FIGDIR/w1_subtopic_trend.png}
\figcap{Per-issue subtopic counts across the week from \texttt{state/metrics.csv}, seven
issues. Daily totals 21, 16, 22, 21, 18, 30 and 20 --- a \textbf{1.9$\times$ range, the
flattest week in the series}, against 2.6$\times$ last week and 5.4$\times$ the week
before. The weekend suppression that dominated the last two rollups is largely absent
here: Sunday 13 September returned 21 records, not the collapse seen on 6 and 12
September, because the Crossref-by-ISSN and Europe PMC legs now carry the weekend
entirely. The 30-record Friday is the 18 September issue, which was built twice --- a
partial morning pass and a full-window re-harvest at 23:0x --- and is a complete day
rather than a deposit anomaly.}

\footnotesize
\textbf{Expanding.} \textit{Exposure assessment \& modelling} \textbf{20 $\rightarrow$
38} is the week's dominant movement and is discussed above. \textit{Other clinical
endpoints} \textbf{4 $\rightarrow$ 8} doubles, driven almost entirely by the 19 September
issue, in which six of twenty records landed in that residual band --- frailty, acute-care
utilisation, renal disease, paediatric respiratory infection, ocular surface and
psoriasis. \textit{Mechanistic toxicology} \textbf{9 $\rightarrow$ 14} and
\textit{Cardiovascular \& metabolic} \textbf{8 $\rightarrow$ 11} both continue the
recovery begun last week. \textit{Respiratory \& allergic} \textbf{12 $\rightarrow$ 13}
holds.

\vspace{1.2mm}
\textbf{Contracting.} \textit{Occupational \& indoor} \textbf{15 $\rightarrow$ 12}
eases for a second consecutive week off the 19-record surge of 5 September, but
\textbf{remains above the pre-surge regular-week maximum of 9 for a third straight week}.
The 5 September rollup declined to call the indoor turn a trend and the 12 September
rollup set a second week as the test; it has now held a third. \textbf{On this watch's own
stated criteria the indoor turn is now a trend and not a fluctuation.}
\textit{Sensing} \textbf{31 $\rightarrow$ 27} and \textit{Burden, policy \& mitigation}
\textbf{14 $\rightarrow$ 13} are mild. \textit{Reproductive \& developmental}
\textbf{6 $\rightarrow$ 5} stays at the bottom of the table, where it has sat every week
of this series.

\vspace{1.2mm}
\textbf{The inflection worth naming.} It is not a subtopic movement at all --- it is that
\textbf{two consecutive weekly rollups set explicit falsifiable tests and both were
answered this week}. \textit{Exposure} was tested against a mid-20s return and overshot
to 38; the indoor turn was tested against a third week and held. A watch that only
reported counts could not have done either. The concrete thing to watch next week is
\textit{Occupational \& indoor}: having now qualified as a trend, the interesting question
is no longer whether it persists but whether any of it moves from chambers and models
into occupied buildings with health outcomes, which so far \textbf{none} of it has.

%% LAYOUT RULE (2026-08-02): architecture donut in the whitespace here, not on its own page.
\vspace{2mm}
\noindent\includegraphics[width=\linewidth]{\FIGDIR/f2a_architecture.png}
\figcap{Study architecture over the pooled week. \textit{Measurement campaign} (29) leads
\textit{Modelling / inventory} (23) and \textit{Review / synthesis} (22), extending the
measurement-ahead-of-modelling ordering restored last week. \textbf{\textit{Metadata only}
falls to 15 and to fourth place}, having been second at 20 a week ago --- the clearest
single-number statement of the screening change. Primary observational epidemiology
(cross-sectional 15, acute 13, cohort 8, ecological 6) totals \textbf{42}, the largest
grouping in the week if read as one. \textbf{One} interventional record appears in the
148, against two last week.}

\clearpage

% -------------------------------------------------------------- pooled analytics
\section*{\sffamily\bfseries\color{Deep}Pooled analytics}
\vspace{-1.5mm}{\color{Teal}\rule{\linewidth}{1.1pt}}\vspace{3mm}

\noindent\includegraphics[width=\linewidth]{\FIGDIR/f1_subtopics.png}
\figcap{Pooled subtopic distribution, \textbf{all ten bands populated}.
\textit{Exposure assessment \& modelling} at 38 and \textit{Sensing} at 27 put the
measurement side at \textbf{65 of 148 (44\%)}. The health side is unusually evenly
spread beneath them --- \textit{Mechanistic toxicology} 14, \textit{Respiratory} 13,
\textit{Burden} 13, \textit{Occupational \& indoor} 12, \textit{Cardiovascular \&
metabolic} 11 --- five clusters inside a four-record band, which is the tightest
health-side distribution in the series and means no single clinical literature dominated
the week.}

\vspace{2mm}
\noindent\includegraphics[width=\linewidth]{\FIGDIR/f2b_endpoint.png}
\figcap{Health endpoint over the pooled week, canonicalised. \textbf{``No health
endpoint'' is 60 of 148 (41\%)}, the aerosol-chemistry, emissions, instrument-performance
and not-applicable records pooled. Among the records that do report a human outcome,
\textbf{respiratory leads at 19}, then other clinical at 10 and cardiovascular at 8. Note
that several labels in this chart are \emph{not} human outcomes despite reading like
metrics --- \textit{air quality / concentration}, \textit{instrument performance},
\textit{sensor bias / accuracy}, \textit{emission flux} --- and the gap between those and
a clinical endpoint is the subject of the heatmap below.}

\vspace{2mm}
\noindent\includegraphics[width=\linewidth]{\FIGDIR/f3_metric_geography.png}
\figcap{Left --- particle metric. \textbf{PM$_{2.5}$ alone is 47 of 148} and unspeciated
PM or emissions terms 62; joint PM$_{2.5}$ + PM$_{10}$ is 18. The tail is the finding:
optical properties 8, \textbf{composition or speciation 6}, \textbf{ultrafine or
size-resolved number 3}, PM$_{10}$ alone 3, bioaerosol 1. Right --- geographic
provenance. \textit{Global / multi-region} at 53 is the largest bar and, as every week,
overstates reach --- it pools genuine multi-country work with reviews that have no field
site and with records whose location is unrecoverable. \textbf{China at 38 is the largest
genuine single-country bar}, ahead of North America at 18 and Europe at 12.}

%% LAYOUT RULE: f6 is short -- it fills the whitespace under f3 rather than trailing the
%% two tall forest panels, where it left page 4 ~55% empty on the first build of this issue.
\vspace{2mm}
\noindent\includegraphics[width=\linewidth]{\FIGDIR/f6_lifecourse.png}
\figcap{Life-course windows across the pooled week; counts are non-exclusive.
\textbf{Working-age adulthood dominates at 18}, as it has every week. The notable reading
is at the young end: \textbf{preconception and in utero is 8} and infancy and early
childhood 6, and unlike last week these are carried by completed epidemiology rather than
by measurement-side records --- 289 non-Hodgkin lymphoma cases in a 3.68-million Spanish
birth series, a zero-inflated gestational PM$_{2.5}$ bias analysis, and a 238-child
e-waste macrophage study. \textbf{Puberty and adolescence remains the weakest window at
2}, as in every rollup in this series.}

\clearpage

\noindent\includegraphics[width=\linewidth]{\FIGDIR/f5_forest.png}
\vspace{1mm}
\noindent\includegraphics[width=\linewidth]{\FIGDIR/f5_forest_2.png}
\figcap{All \textbf{47} pooled effect estimates on a common logarithmic axis, split across
two panels. \textbf{These are not mutually comparable and are not meta-analysed}:
contrasts include per-IQR, per-10\,$\mu$g\,m$^{-3}$, per-1\,$\mu$g\,m$^{-3}$, per-SD,
per-exposure-category and mixture-index units, on OR, HR, RR, IRR and ratio scales.
\textbf{Heterogeneity here is structural, not statistical}, so no $I^2$ is reported and
pooling these would be a category error. Four things are readable. First, precision
improved: \textbf{42 of 47 report an interval (89\%)}, and \textbf{four of the five bare
points come from one paper} --- the wildfire versus non-wildfire PM$_{2.5}$ hospitalisation
analysis, which publishes eight headline risks without a single interval. Second,
\textbf{the largest effects are again the least precise}: the widest interval on the axis
belongs to the largest point estimate, \textbf{OR 4.09 (1.06--15.79)} for microplastic
detection in primary lung malignancy, an interval spanning fifteenfold and resting on a
handful of cases. Third, \textbf{the tightest interval belongs to one of the smallest
effects} --- \textbf{HR 1.0171 (1.0149--1.0192)} for all-cause death in lung cancer per
1\,$\mu$g\,m$^{-3}$ of non-wildfire PM$_{2.5}$ --- which is the characteristic shape of
this literature: population-scale precision on effects of a few percent. Fourth, several
estimates are \emph{co-pollutant} rather than PM (O$_3$, NO$_2$, SO$_2$, CO) and are
plotted because the papers report them alongside PM. Read the labels before the positions.}

\vspace{2mm}
\noindent\includegraphics[width=\linewidth]{\FIGDIR/f4_heatmap.png}
\figcap{Subtopic $\times$ architecture, pooled. \textbf{The densest cell in the week is
\textit{Exposure assessment \& modelling} $\times$ \textit{Measurement campaign} at 17},
displacing last week's \textit{Sensing} $\times$ \textit{Metadata only} (which has fallen
from 10 to 4). That single swap is the clearest structural evidence of the screening
change. The standing finding survives an eighth week and is now measured on a larger base:
of the \textbf{65 records} in the \textit{Sensing} and \textit{Exposure} rows,
\textbf{4 carry any human health endpoint, and only 3 of those observe one} --- the fourth
is a modelled inhalation cancer risk, not a measured outcome. Instrumentation and
epidemiology continue to be written by different people in different journals.}

\vspace{4mm}

% -------------------------------------------------------------- venues + trials
\section*{\sffamily\bfseries\color{Deep}Venues, trial watch, and what the pipeline did}
\vspace{-1.5mm}{\color{Teal}\rule{\linewidth}{1.1pt}}\vspace{2.5mm}

\footnotesize
\textbf{Venues.} \textit{Environ Sci Technol} \textbf{11 of 148} leads, followed by
\textit{Environ Pollut} (9), \textit{Front Public Health} (8), \textit{Atmos Chem Phys}
(6) and \textit{Atmosphere} (5), then a four-way tie at 4 between \textit{Atmos Meas
Tech}, \textit{Aerosol Air Qual Res}, \textit{Environ Monit Assess} and
\textit{Zhonghua Liu Xing Bing Xue Za Zhi}. Concentration is low --- the leading venue
holds 7\% of the week --- and the distribution is a harvesting artefact as much as an
editorial one: titles on the Crossref by-ISSN list enter at their own publication rate,
those reachable only through PubMed enter at PubMed's. The presence of \textit{AMT} and
\textit{ACP} at all is entirely due to the by-ISSN leg, since PubMed does not index them.
\textbf{Preprints supplied 2 records}, down from 8, both labelled in the register.

\vspace{1.5mm}
\textbf{Recurring authors.} Unchanged and still unresolvable. Surname frequency is led by
Zhang, Li, Huang and Chen at two or three each, which is a naming artefact of a corpus
that is 26\% Chinese-affiliated and not a research-group signal. \textbf{No repeat-group
cluster is claimed for this week.} Adding a first-author affiliation field remains the
open schema item that would make this section mean anything; it has been on the list since
3 September and is now the oldest unaddressed gap in the series.

\vspace{1.5mm}
\textbf{Trial watch --- five in-window records and, for the first time, nothing new to
state.} ClinicalTrials.gov returned \textbf{5 PM or air-pollution registrations or updates
across 13--19 September}, and \textbf{all five were already in
\texttt{state/trials.json}}, captured by the daily sweeps during the week. That is a
first: on every previous weekly reconciliation the weekly sweep found records the daily
leg had missed. \textbf{The daily trial leg is now ahead of the weekly sweep rather than
behind it}, which is the result the daily sweeps were added to produce. The four topical
records are \textbf{NCT07540000} (INTERLUNG, Johns Hopkins, Bhaktapur, Nepal, $n=800$, a
PHASE4 community-health-worker respiratory intervention with a spirometric primary
endpoint --- still the strongest trial in the accumulated set), \textbf{NCT07793123}
(Robertet SA, oral antioxidant versus placebo on skin radiance, $n=60$),
\textbf{NCT07713264} (VISIONS, Stanford, outdoor-worker education, $n=50$) and
\textbf{NCT07817836} (Pavia, completed observational, DNA-methylation biological age,
$n=151$). The fifth, \textbf{NCT02820402} (Project Viva), is a 1999 birth cohort whose
air-pollution work is secondary and is deliberately not stored. \textbf{Three of the four
topical records measure no air-pollution variable at all} --- exposure is inferred from
urban residence or occupation --- which is the fourth consecutive week in which the
registry's air-pollution trials are mostly trials about something else.
\texttt{trials.json} now holds \textbf{12 trials across 14 windows}.

\vspace{1.5mm}
\textbf{What the harvesting legs did.} \textbf{Crossref-by-ISSN was again the workhorse},
highest-yield leg on every day it ran and the only route to the instrumentation journals.
\textbf{Europe PMC} was the second leg and the one that made the weekend survivable ---
Sunday 13 September returned 21 records where the two previous Sundays collapsed.
\textbf{PubMed} ran as two separate queries on the health and instrumentation axes all
week; \textbf{the instrumentation query returned zero on every single day}, which after
seven consecutive days is no longer a sampling observation but a confirmation that PubMed
structurally does not index this literature. \textbf{OpenAlex returned HTTP 429 every
day}, unchanged since 28 July. \textbf{arXiv returned HTTP 406 every day} and contributed
no record to any issue this week --- third consecutive week of that, and the query is
overdue for inspection. \textbf{Consensus returned a nil every day} for the fourth
consecutive week, its hits consistently the right literature with Crossref
\texttt{created} dates years outside the window. \textbf{One new abstract-recovery route
was tested and failed}: Semantic Scholar, queried per DOI on 19 September, returned
\textbf{404 on all five} same-day deposits, having not yet indexed them. Together with the
OpenAlex per-work failure logged on 18 September, that closes the two cheapest remaining
recovery options, and leaves the deferred-retry pass as the only mechanism that can work.

\vspace{2mm}
% -------------------------------------------------------------- gaps
\section*{\sffamily\bfseries\color{Deep}Open gaps}
\vspace{-1.5mm}{\color{Teal}\rule{\linewidth}{1.1pt}}\vspace{2.5mm}

\footnotesize
\begin{itemize}
\item \textbf{The abstract-retry pass is still not built, and is now the only thing that
can work.} The 10\% metadata-only rate is an improvement produced by \emph{rejecting}
retro-indexed records, not by \emph{recovering} withheld ones. Two recovery routes were
tested and failed this week and last --- OpenAlex per-work, and Semantic Scholar per DOI
--- both for the same reason: they cannot return an abstract the publisher never deposited
or that no index has yet ingested. \textbf{The only route left is time}: re-query Crossref,
Europe PMC and PubMed for the standing held list a week after first sight. That list now
holds \textbf{18 records}.

\item \textbf{Four of 65 measurement-side records carry a human health endpoint, and only
three observe one.} Eighth consecutive week, now on the largest single-week base yet
recorded. This remains the most durable structural finding in the series.

\item \textbf{Ultrafine and size-resolved number: 3 of 148.} Unchanged as a share from
last week. The one instrumentation record in the 19 September issue that would speak
directly to this --- a sub-10\,nm unipolar aerosol charger in \textit{J Aerosol Sci} ---
shipped without an abstract, which is the gap and the pipeline failure landing on the
same record.

\item \textbf{One interventional record in 148}, and it does not test a monitoring
technology. The \textit{Occupational \& indoor} cluster now qualifies as a trend on this
watch's own criteria, and \textbf{none} of it has yet been evaluated against health
outcomes in occupied buildings rather than against chambers and models.

\item \textbf{Pipeline defects carried forward}: no first-author affiliation field (oldest
open item); the \texttt{Cancer}/\texttt{Oncologic} split in \texttt{ENDPOINT\_CANON}; a
6-element \texttt{LIFECOURSE} array in the 2026-08-03 corpus where the schema requires 5
--- outside this window but it \textbf{will} corrupt the September monthly; and
\texttt{.git} lock and temp debris that the delete-restricted mount cannot remove, wanting
manual cleanup on a local checkout. \textbf{Resolved this week}: \texttt{metrics.csv} has
been verified unquoted-safe for thirteen consecutive issues (428 rows, 0 malformed), so
the ``known defect'' note in the task specification is stale and should be struck.
\end{itemize}

\clearpage

% -------------------------------------------------------------- digest
\section*{\sffamily\bfseries\color{Deep}Paper-level digest --- the week by subtopic}
\vspace{-1.5mm}{\color{Teal}\rule{\linewidth}{1.1pt}}\vspace{1mm}

{\fontsize{7.4}{9}\selectfont\color{Slate}Every record that appeared in a daily issue
keeps that issue's critical summary verbatim; nothing here is re-written or re-assessed.
Records surfaced without an abstract are shown as such. Ordered by subtopic, then by the
date of the issue that first carried the record.\par}

\input{digest_body}

% -------------------------------------------------------------- register
\section*{\sffamily\bfseries\color{Deep}Pooled corpus register}
\vspace{-1.5mm}{\color{Teal}\rule{\linewidth}{1.1pt}}\vspace{2.5mm}

{\fontsize{6.7}{8.4}\selectfont
\setlength{\tabcolsep}{2.4pt}
\begin{longtable}{@{}>{\RaggedRight}p{25mm}>{\RaggedRight}p{41mm}>{\RaggedRight}p{27mm}%
>{\RaggedRight}p{23mm}>{\RaggedRight}p{18mm}>{\RaggedRight}p{15mm}@{}}
\toprule
\textsf{\textbf{First author}} & \textsf{\textbf{Focus}} & \textsf{\textbf{Journal}} &
\textsf{\textbf{Design}} & \textsf{\textbf{Metric}} & \textsf{\textbf{Region}}\\
\midrule
\endfirsthead
\toprule
\textsf{\textbf{First author}} & \textsf{\textbf{Focus}} & \textsf{\textbf{Journal}} &
\textsf{\textbf{Design}} & \textsf{\textbf{Metric}} & \textsf{\textbf{Region}}\\
\midrule
\endhead
\multicolumn{6}{@{}l@{}}{\cellcolor{Mist}\textsf{\textbf{\textcolor{Deep}{Neuro / mental health}}}}\\[0.6mm]
\href{https://doi.org/10.1001/jamanetworkopen.2026.25015}{Dales} & Wildfire-sourced PM$_{2.5}$ and ED visits for migraine and other primary headache & \textit{JAMA Netw Open} & Case-crossover & PM$_{2.5}$ (wildfire) & Canada\\
\href{https://doi.org/10.1093/ageing/afag215}{Glover} & Umbrella meta-review of urban environmental exposures and cognitive health & \textit{Age Ageing} & Umbrella review / meta-analysis & PM$_{2.5}$ & Global\\
\href{https://doi.org/10.1016/j.nsa.2026.107019}{Goossens} & In utero air pollution exposure and pre-/postnatal brain development & \textit{Neurosci Appl} & Narrative review & PM$_{2.5}$ / PM$_{10}$ / PM$_{\mathrm{coarse}}$ & Global\\
\href{https://doi.org/10.1016/j.isci.2026.116617}{Ma} & PM$_{2.5}$-induced neurotoxicity and depression: exposure risks and mechanisms & \textit{iScience} & Narrative review & PM$_{2.5}$ + components & Global\\
\href{https://doi.org/10.3760/cma.j.cn112150-20250804-00752}{Pan} & Maternal PM exposure and offspring neurodevelopment: effects and mechanisms & \textit{Zhonghua Yu Fang Yi Xue Za Zhi} & Narrative review & PM (bulk) & China\\
\addlinespace[0.9mm]
\multicolumn{6}{@{}l@{}}{\cellcolor{Mist}\textsf{\textbf{\textcolor{Teal}{Cardiovascular \& metabolic}}}}\\[0.6mm]
\href{https://doi.org/10.3760/cma.j.cn112338-20251120-00843}{Huang} & PM$_{2.5}$ and dyslipidaemia in rural Chinese children; dietary effect modification & \textit{Zhonghua Liu Xing Bing Xue Za Zhi} & Cross-sectional (multistage) & PM$_{2.5}$ & China\\
\href{https://doi.org/10.32394/pe/218966}{Kazmierska} & Air pollution and cardiovascular disease: review of current data & \textit{Przegl Epidemiol} & Narrative review & PM$_{2.5}$ / PM$_{10}$ & Global\\
\href{https://doi.org/10.1007/s00125-026-06814-2}{Xi} & Air pollution, built environment and temperature vs. prediabetes phenotypes & \textit{Diabetologia} & Prospective cohort & PM$_{2.5}$ & Germany\\
\href{https://doi.org/10.1161/JAHA.125.046917}{Yang} & Short-term PM$_{2.5}$ and hospitalisation/ED use in atrial fibrillation & \textit{J Am Heart Assoc} & Case-crossover & PM$_{2.5}$ (wildfire-influenced) & USA\\
\addlinespace[0.9mm]
\multicolumn{6}{@{}l@{}}{\cellcolor{Mist}\textsf{\textbf{\textcolor{Sky}{Reproductive \& developmental}}}}\\[0.6mm]
\href{https://doi.org/10.1016/j.jes.2025.10.043}{Chen} & Prenatal air pollution, fetal growth trajectory phenotypes and child neurodevelopment & \textit{J Environ Sci (China)} & Prospective cohort & PM$_{2.5}$ / PM$_{10}$ & China\\
\href{https://doi.org/10.1016/j.ecoenv.2026.120566}{Huang} & Joint prenatal bisphenol + air-pollutant mixture exposure and preschool IQ & \textit{Ecotoxicol Environ Saf} & Prospective cohort (mixtures) & PM$_{2.5}$ / PM$_{10}$ & China\\
\href{https://doi.org/10.1007/s40572-026-00552-8}{Peterson} & Ambient air pollution, PAHs and pubertal development: critical review & \textit{Curr Environ Health Rep} & Critical review & PM$_{2.5}$ components & Global\\
\href{https://doi.org/10.1016/j.envres.2026.125313}{Xie} & Long-term PM$_{2.5}$ and timing of pubertal development in boys and girls & \textit{Environ Res} & Prospective cohort & PM$_{2.5}$ & China\\
\addlinespace[0.9mm]
\multicolumn{6}{@{}l@{}}{\cellcolor{Mist}\textsf{\textbf{\textcolor{Sage}{Respiratory \& allergic}}}}\\[0.6mm]
\href{https://doi.org/10.1007/s00484-026-03289-z}{Bayar Muluk} & Climate change and otolaryngology: airway disease impacts and emission reduction & \textit{Int J Biometeorol} & Narrative review & PM$_{2.5}$ & Global\\
\href{https://doi.org/10.4103/mgr.MEDGASRES-D-25-00164}{Liao} & Short-term air pollution, respiratory symptoms and lung function in asthma & \textit{Med Gas Res} & Panel study & PM$_{2.5}$ / PM$_{10}$ & China\\
\href{https://doi.org/10.3390/epidemiologia7040097}{Sanchez de Prada} & Ambient PM and asthma-coded admissions in children 0-3 y, Spain & \textit{Epidemiologia} & Ecological panel & PM$_{2.5}$ / PM$_{10}$ & Spain\\
\href{https://doi.org/10.1016/j.jhazmat.2026.143076}{Yang} & Allergenic vs. toxic risk of airborne seed hairs compared with ambient PM$_{2.5}$ & \textit{J Hazard Mater} & Field measurement & PM$_{2.5}$ + bioaerosol & China\\
\addlinespace[0.9mm]
\multicolumn{6}{@{}l@{}}{\cellcolor{Mist}\textsf{\textbf{\textcolor{Violet}{Mechanistic toxicology}}}}\\[0.6mm]
\href{https://doi.org/10.1016/j.jes.2025.12.074}{Ding} & PM$_{2.5}$ promotes vascular calcification via miR-27b-5p/HES1 signalling & \textit{J Environ Sci (China)} & In vitro + organoid + murine & PM$_{2.5}$ & China\\
\href{https://doi.org/10.1007/s00360-026-01698-5}{Maraschi} & Osmoregulatory disruption and metal bioaccumulation from settleable atmospheric PM & \textit{J Comp Physiol B} & Controlled exposure (ecotox) & Settleable APM & Brazil\\
\href{https://doi.org/10.1016/j.jes.2025.11.018}{Wu} & Macrophage NLRP3 inflammasome in real-ambient PM-induced hepatic glucose dysregulation & \textit{J Environ Sci (China)} & Real-ambient murine + co-culture & Real-ambient PM & China\\
\addlinespace[0.9mm]
\multicolumn{6}{@{}l@{}}{\cellcolor{Mist}\textsf{\textbf{\textcolor{Clay}{Occupational \& indoor}}}}\\[0.6mm]
\href{https://doi.org/10.3389/fmedt.2026.1796828}{Chen} & Nested suction instrument to cut intraoperative surgical-smoke PM$_{2.5}$ & \textit{Front Med Technol} & Bench / device experiment & PM$_{2.5}$ (surgical smoke) & China\\
\href{https://doi.org/10.1093/annweh/wxag055}{Cole} & Determinants of respirable crystalline silica exposure in Australian tunnelling & \textit{Ann Work Expo Health} & Observational exposure study & Respirable dust / RCS & Australia\\
\href{https://doi.org/10.1016/j.jhazmat.2026.143081}{Ding} & Stage-resolved hair Zn-isotope dynamics as a firefighting combustion-exposure tracer & \textit{J Hazard Mater} & Repeated-measures biomarker & Combustion-derived PM & China\\
\href{https://doi.org/10.3390/jox16040119}{Ramsden} & Xenobiotic hazards in aircraft cabin air & \textit{J Xenobiot} & Review + risk economics & Ultrafine particles & Global\\
\addlinespace[0.9mm]
\multicolumn{6}{@{}l@{}}{\cellcolor{Mist}\textsf{\textbf{\textcolor{Amber}{Sensing, forecasting \& instrumentation}}}}\\[0.6mm]
\href{https://doi.org/10.1016/j.jenvman.2026.130515}{Shi} & Adaptive spatiotemporal graph transformer for multi-site multi-step PM$_{2.5}$ forecasting & \textit{J Environ Manage} & Deep-learning model & PM$_{2.5}$ & China / India\\
\href{https://doi.org/10.3390/bios16070380}{Tsuchiya} & Seasonal and indoor-outdoor characteristics of fluorescent aerosol particles in offices & \textit{Biosensors} & Real-time sensor campaign & Bioaerosol / coarse & Japan\\
\addlinespace[0.9mm]
\multicolumn{6}{@{}l@{}}{\cellcolor{Mist}\textsf{\textbf{\textcolor{Amber}{Exposure assessment \& modelling}}}}\\[0.6mm]
\href{https://doi.org/10.1029/2025GH001621}{Solorzano-Araque} & Multiscale exposure modelling: LOTOS-EUROS + mobility surveys + particle toxicity & \textit{GeoHealth} & CTM + mobility integration & PM$_{2.5}$ / PM$_{10}$ & Colombia\\
\addlinespace[0.9mm]
\multicolumn{6}{@{}l@{}}{\cellcolor{Mist}\textsf{\textbf{\textcolor{Coral}{Burden, policy \& mitigation}}}}\\[0.6mm]
\href{https://doi.org/10.1002/hsr2.72886}{Karamzad} & Ischaemic heart disease burden in MENA 1990-2021; PM attributable fraction 28.2\% & \textit{Health Sci Rep} & GBD secondary analysis & Ambient PM & MENA\\
\href{https://doi.org/10.1371/journal.pone.0347989}{Rees} & HEPA filtration to reduce respiratory infection in care homes & \textit{PLoS One} & RCT process evaluation & Filtered PM & UK\\
\href{https://doi.org/10.1016/j.jes.2026.01.009}{Shan} & Policy-driven emission reductions and post-pandemic rebound, high-resolution inventories & \textit{J Environ Sci (China)} & Emission inventory & PM emissions & China\\
\addlinespace[0.9mm]
\multicolumn{6}{@{}l@{}}{\cellcolor{Mist}\textsf{\textbf{\textcolor{Slate}{Other clinical endpoints}}}}\\[0.6mm]
\href{https://doi.org/10.3390/dj14070440}{Lupita} & PM$_{2.5}$, preventive-care capacity and caries/edentulism burden across the EU & \textit{Dent J} & Ecological panel & PM$_{2.5}$ & EU-27\\
\href{https://doi.org/10.1016/j.jdin.2026.06.008}{Park} & Long-term PM exposure and risk of androgenetic alopecia & \textit{JAAD Int} & Nationwide cohort & PM$_{2.5}$ / PM$_{10}$ & South Korea\\
\href{https://doi.org/10.3171/2026.3.PEDS2615}{Sun} & Environmental exposure and area deprivation vs. paediatric hydrocephalus outcomes & \textit{J Neurosurg Pediatr} & Registry cohort + geospatial & PM$_{2.5}$ (satellite-derived) & USA\\
\addlinespace[0.9mm]
\bottomrule
\end{longtable}}


\vspace{4mm}
\begin{keybox}
{\sffamily\bfseries\footnotesize\color{Deep}Provenance and method}\\[1.2mm]
{\fontsize{7.2}{8.8}\selectfont
Period \textbf{13--19 September 2026}, Sunday to Saturday, \textbf{seven daily issues},
\textbf{148 records}, \textbf{148 unique DOIs}, no record without a DOI. This rollup is
\textbf{aggregation only}: it re-reads \texttt{state/corpus/*.json},
\texttt{state/metrics.csv} and \texttt{state/trials.json} over the period and
\textbf{queries no literature API}, so every record and every number here was gated by
\texttt{check\_dois.py} at daily-issue time and nothing has been re-screened or
re-assessed. The only live query made for this issue was the ClinicalTrials.gov
reconciliation sweep reported above, which is state maintenance and contributes no record
to the corpus. Tier mix across the week: \textbf{35 A, 71 B, 42 C}. The 19 September issue
was a recovery build for a scheduled run that did not fire; it covers its full day and the
week therefore has no gap. Effect estimates are transcribed verbatim from the daily issues
and use heterogeneous exposure contrasts --- see the forest-plot caption --- and are
\textbf{not} meta-analysed. Abstracts remain the copyright of their respective
publishers.\par}
\end{keybox}

\end{document}
